\documentclass[aspectratio=169,11pt]{beamer}
\usepackage[T1]{fontenc}
\usepackage{lmodern}
\usepackage{listings}
\usepackage{booktabs}
\definecolor{navy}{RGB}{25,48,76}
\definecolor{teal}{RGB}{0,112,120}
\definecolor{codebg}{RGB}{246,247,249}
\setbeamercolor{structure}{fg=teal}
\setbeamercolor{frametitle}{fg=navy}
\setbeamercolor{title}{fg=navy}
\setbeamercolor{normal text}{fg=navy,bg=white}
\setbeamertemplate{navigation symbols}{}
\setbeamertemplate{footline}{\hfill\usebeamercolor[fg]{structure}\insertframenumber\,/\,\inserttotalframenumber\hspace{5mm}\vspace{3mm}}
\setbeamertemplate{itemize items}[circle]
\setbeamersize{text margin left=9mm,text margin right=9mm}
\lstset{language=Python,basicstyle=\ttfamily\small,keywordstyle=\color{teal}\bfseries,commentstyle=\color{gray},stringstyle=\color{teal},backgroundcolor=\color{codebg},frame=none,showstringspaces=false,columns=fullflexible,keepspaces=true,breaklines=true,xleftmargin=2mm,xrightmargin=2mm,aboveskip=8pt,belowskip=8pt}
\newcommand{\code}[1]{\texttt{#1}}
\usepackage{amsmath}
\title{Testing Python with pytest}
\subtitle{Assertions, fixtures, and tests for language tools}
\author{CS413}
\date{October 8, 2026}
\begin{document}
\begin{frame}\titlepage\end{frame}
\begin{frame}[fragile]{Testing as an executable specification}
\begin{itemize}
\item A test describes one observable behavior and checks it automatically.
\item Example: looking up a missing variable should raise a particular exception.
\item A failing test gives a counterexample to the claimed behavior.
\item Passing tests provide evidence for the chosen cases, not a proof for all inputs.
\end{itemize}
\end{frame}

\begin{frame}[fragile]{Today's goals}
\begin{itemize}
\item Write and run tests with ordinary Python assertions.
\item Choose boundary cases and express them with parametrization.
\item Use fixtures to control inputs, resources, and side effects.
\item Add regression tests for a small language implementation.
\end{itemize}
\end{frame}

\begin{frame}[fragile]{A small project}
\begin{lstlisting}
pytest_demo/
    stack.py
    test_stack.py
    test_resources.py
    conftest.py
    expr.py
    test_expr.py
\end{lstlisting}
Run commands from \code{examples/pytest\_demo/}.
\begin{lstlisting}
python3 -m venv .venv
source .venv/bin/activate  # POSIX shell
python -m pip install pytest
python -m pytest -q
\end{lstlisting}

\end{frame}

\begin{frame}[fragile]{The behavior we want to test}
\begin{lstlisting}
# stack.py
class Stack:
    def __init__(self):
        self._items = []

    def push(self, value):
        self._items.append(value)

    def pop(self):
        if not self._items:
            raise IndexError("empty stack")
        return self._items.pop()
\end{lstlisting}
Contract: last in, first out. An empty pop raises \code{IndexError}.
\end{frame}

\begin{frame}[fragile]{A first test}
\begin{lstlisting}
# test_stack.py
from stack import Stack

def test_pop_returns_last_pushed():
    stack = Stack()           # arrange
    stack.push("x")
    stack.push("y")
    result = stack.pop()      # act
    assert result == "y"      # assert
\end{lstlisting}
The test name describes the behavior. A failing assertion fails the test.
\end{frame}

\begin{frame}[fragile]{Discovery and execution}
\begin{itemize}
\item By default, pytest searches for \code{test\_*.py} and \code{*\_test.py}.
\item Test functions start with \code{test\_}. Plain functions are enough.
\item pytest imports test modules, collects tests, then runs them.
\item Keep demonstrations out of module-level code to avoid import side effects.
\end{itemize}\begin{lstlisting}
python -m pytest --collect-only -q
\end{lstlisting}

\end{frame}

\begin{frame}[fragile]{Reading a failure}
Suppose \code{pop()} incorrectly uses \code{self.\_items.pop(0)}.
\begin{lstlisting}
>       assert result == "y"
E       AssertionError: assert 'x' == 'y'

FAILED test_stack.py::test_pop_returns_last_pushed
\end{lstlisting}
\begin{itemize}
\item The node ID identifies the file and test function.
\item Read the actual and expected values, then inspect the contract.
\item An exception during fixture setup is reported as an error.
\end{itemize}
\end{frame}

\begin{frame}[fragile]{Running a focused set of tests}
\begin{lstlisting}
python -m pytest -q
python -m pytest test_stack.py -v
python -m pytest test_stack.py::test_empty_pop
python -m pytest -k "pop and not empty"
python -m pytest -x
python -m pytest --lf
\end{lstlisting}
\begin{itemize}
\item \code{-k} selects by a name expression. \code{-x} stops on the first failure.
\item \code{--lf} reruns previously failing tests using the local cache.
\item During debugging, run a small selection. Before finishing, run the suite.
\end{itemize}
\end{frame}

\begin{frame}[fragile]{Assertions and the test oracle}
\begin{lstlisting}
def test_pop_removes_an_item():
    stack = Stack()
    stack.push("x")
    stack.push("y")
    assert stack.pop() == "y"
    assert stack.pop() == "x"
\end{lstlisting}
\begin{itemize}
\item The oracle is how the test decides the expected result.
\item Choose expectations from the specification, independently of the implementation.
\item Checking only that a result exists can miss an incorrect result.
\end{itemize}
\end{frame}

\begin{frame}[fragile]{Expected exceptions}
\begin{lstlisting}
import pytest

def test_empty_pop():
    stack = Stack()
    with pytest.raises(IndexError, match="empty stack"):
        stack.pop()
\end{lstlisting}
\begin{itemize}
\item No exception, or a different exception type, fails this test.
\item \code{match} is a regular expression search on the exception message.
\item Keep the block narrow. Put assertions about later behavior after it.
\end{itemize}
\end{frame}

\begin{frame}[fragile]{Several cases, one test}
\begin{lstlisting}
@pytest.mark.parametrize("values", [
    [],
    [7],
    [2, 5, 2],
], ids=["empty", "singleton", "duplicates"])
def test_lifo(values):
    stack = Stack()
    for value in values:
        stack.push(value)
    actual = [stack.pop() for _ in values]
    assert actual == list(reversed(values))
\end{lstlisting}
pytest collects a separate test for each case, with its own result.
\end{frame}

\begin{frame}[fragile]{Choosing cases that expose mistakes}
\begin{itemize}
\item Empty input exercises a boundary or error condition.
\item A singleton checks the smallest successful operation.
\item Distinct values reveal ordering mistakes. Duplicates reveal lost elements.
\item After pushing and popping everything, another pop should fail.
\item A regression test reproduces a specific bug before the fix and passes after it.
\end{itemize}
\end{frame}

\begin{frame}[fragile]{Fixtures supply test inputs}
\begin{lstlisting}
# conftest.py
import pytest
from stack import Stack

@pytest.fixture
def stack():
    return Stack()

# test_stack.py
def test_fixture_stack(stack):
    stack.push(7)
    assert stack.pop() == 7
\end{lstlisting}
pytest resolves the argument name \code{stack} and supplies the fixture value.
\end{frame}

\begin{frame}[fragile]{Fixture lifetime and sharing}
\begin{itemize}
\item The default scope is \code{function}: each test gets a fresh stack.
\item A fixture may request another fixture through its own arguments.
\item Put shared fixtures in \code{conftest.py}. Tests do not import that file.
\item Wider scopes include \code{class}, \code{module}, \code{package}, and \code{session}.
\item A mutable object shared across tests can make results depend on execution order.
\end{itemize}
\end{frame}

\begin{frame}[fragile]{Cleanup with a yield fixture}
\begin{lstlisting}
@pytest.fixture
def output_file(tmp_path):
    stream = (tmp_path / "out.txt").open("w+")
    yield stream
    stream.close()

def test_output_file(output_file):
    output_file.write("hello")
    output_file.seek(0)
    assert output_file.read() == "hello"
\end{lstlisting}
Setup runs before \code{yield}. Cleanup runs afterward, even if the test fails.

If setup fails before reaching \code{yield}, that trailing cleanup does not run.
\end{frame}

\begin{frame}[fragile]{Temporary files with tmp\_path}
\begin{lstlisting}
def test_save_names(tmp_path):
    path = tmp_path / "names.txt"
    path.write_text("Ada\nGrace\n", encoding="utf-8")
    names = path.read_text(encoding="utf-8").splitlines()
    assert names == ["Ada", "Grace"]
\end{lstlisting}
\begin{itemize}
\item \code{tmp\_path} provides a separate \code{pathlib.Path} directory per test.
\item Use it for file-based tests instead of a fixed filename in the project.
\item This small example illustrates the fixture. Application tests should call the application's file-handling code.
\end{itemize}
\end{frame}

\begin{frame}[fragile]{Printed output with capsys}
\begin{lstlisting}
def announce(name):
    print(f"Hello, {name}!")

def test_announce(capsys):
    announce("Ada")
    captured = capsys.readouterr()
    assert captured.out == "Hello, Ada!\n"
    assert captured.err == ""
\end{lstlisting}
\code{capsys} captures Python standard output and standard error.

The newline is part of the expected output.
\end{frame}

\begin{frame}[fragile]{Environment control with monkeypatch}
\begin{lstlisting}
# stack.py
import os

def greeting():
    return "Hello, " + os.getenv("USER_NAME", "guest")

# test_resources.py
from stack import greeting

def test_greeting(monkeypatch):
    monkeypatch.setenv("USER_NAME", "Ada")
    assert greeting() == "Hello, Ada"
\end{lstlisting}
pytest restores the environment change after the test.
\end{frame}

\begin{frame}[fragile]{Replacing dependencies}
\begin{itemize}
\item Unit tests can replace a clock, network call, or random source with a controlled substitute.
\item \code{monkeypatch.setattr} temporarily replaces an attribute.
\item Patch the name where the code under test looks it up.
\item Keep integration tests that exercise the real boundary too.
\item A substitute should preserve the relevant interface and behavior. Otherwise a passing test may be misleading.
\end{itemize}
\end{frame}

\begin{frame}[fragile]{Floating-point results}
\begin{lstlisting}
import pytest

def test_decimal_sum():
    assert 0.1 + 0.2 == pytest.approx(0.3, abs=1e-12,
                                   rel=0)
\end{lstlisting}
\begin{itemize}
\item Binary floating-point arithmetic can introduce rounding error.
\item Choose a tolerance that fits the numerical contract.
\item Use exact equality for exact results, including integer counts and sets.
\end{itemize}
\end{frame}

\begin{frame}[fragile]{Skip and expected failure}
\begin{lstlisting}
@pytest.mark.skip(reason="requires the course server")
def test_server():
    ...

@pytest.mark.xfail(strict=True, reason="issue 42: known bug")
def test_known_bug():
    assert buggy_function() == expected
\end{lstlisting}
\begin{itemize}
\item Illustrative examples: the dependency and bug must exist in a real suite.
\item \code{skip} does not run the test. \code{xfail} normally runs it expecting failure.
\item With \code{strict=True}, an unexpected pass fails the suite so you revisit the marker.
\end{itemize}
\end{frame}

\begin{frame}[fragile]{A small language example}
Our standalone demo represents expressions with tuples:
\begin{lstlisting}
("var", "x")
("app", function_expr, argument_expr)
("lam", "x", body_expr)
\end{lstlisting}
The free-variable rules are:
\[
\begin{aligned}
FV(x) &= \{x\} \\
FV(e_1\ e_2) &= FV(e_1) \cup FV(e_2) \\
FV(\lambda x.e) &= FV(e) \setminus \{x\}
\end{aligned}
\]
The same expectations can test a visitor-based implementation.
\end{frame}

\begin{frame}[fragile]{Exercise: free-variable tests}
\begin{lstlisting}
from expr import free_vars

@pytest.mark.parametrize("expr, expected", [
    (("var", "x"), {"x"}),
    (("lam", "x", ("var", "x")), set()),
    # Add an application and a nested lambda.
])
def test_free_vars(expr, expected):
    assert free_vars(expr) == expected
\end{lstlisting}
\begin{itemize}
\item Write a case where a variable remains free inside a lambda.
\item Which case catches a lambda rule that forgets to remove its binder?
\item Work for five minutes, then compare cases with a neighbor.
\end{itemize}
\end{frame}

\begin{frame}[fragile]{Exercise discussion}
\begin{lstlisting}
# Application combines both sets.
(("app", ("var", "x"), ("var", "y")), {"x", "y"})

# Binding x leaves y free.
(("lam", "x", ("var", "y")), {"y"})

# Both occurrences of x are bound by the outer lambda.
(("lam", "x", ("app", ("var", "x"),
    ("lam", "y", ("var", "x")))), set())
\end{lstlisting}
A test of \code{("lam", "x", ("var", "x"))} catches the missing subtraction.

Set equality checks membership without depending on display order.
\end{frame}

\begin{frame}[fragile]{A reliable testing workflow}
\begin{itemize}
\item Describe the intended behavior with a small example.
\item Write the test and observe the expected failure.
\item Change the implementation and rerun the focused test.
\item Run the whole suite before committing the change.
\item Keep tests independent of execution order, external state, and wall-clock timing.
\end{itemize}
\end{frame}

\begin{frame}[fragile]{What a passing suite tells us}
\begin{itemize}
\item The implementation satisfies the assertions on the inputs exercised.
\item Unasserted behavior and untested cases can still contain bugs.
\item Code coverage measures execution, not whether the assertions are strong.
\item A useful check: deliberately introduce a plausible bug and see whether a test catches it.
\end{itemize}
\end{frame}

\begin{frame}[fragile]{Reference and practice}
\begin{itemize}
\item Demo: \code{examples/pytest\_demo/}. Run \code{python -m pytest -q} there.
\item Practice: add a stack test that interleaves pushes and pops.
\item Practice: add a nested-lambda case with both a bound and a free variable.
\end{itemize}\medskip
Official pytest documentation:
\begin{itemize}
\item \href{https://docs.pytest.org/en/stable/getting-started.html}{Getting started and discovery}
\item \href{https://docs.pytest.org/en/stable/how-to/assert.html}{Assertions} and \href{https://docs.pytest.org/en/stable/how-to/parametrize.html}{parametrization}
\item \href{https://docs.pytest.org/en/stable/how-to/fixtures.html}{Fixtures} and \href{https://docs.pytest.org/en/stable/how-to/monkeypatch.html}{monkeypatch}
\end{itemize}
\end{frame}
\end{document}
